\section{The plane: an audited finite proof and its constructive scope}
\label{sec:plane-new}

The repository already contains a substantial planar development
\cite{prior-planar}.  The
unchanged source is included as \texttt{reused/isoperimetric\_ipra.tex},
and its complete compiled paper accompanies this report.  It was retrieved
at ProveIt commit
\texttt{22afac4c9db7a0e060c081cc11db0dc96fe59e93}; the source itself
records the older foundation snapshot on which it was originally based.
The discussion below gives the decisive finite proof and the audited
equality estimate, so the dependence is mathematical and explicit.
The original paper supplies the detailed positive-radical coding,
lower-real perimeter construction, and the remaining coding details.
Its status is a human proof with a Lean outline, not a completed
internal Lean derivation.

\subsection{A theorem about rational polygons}
For a cyclic list of rational points \(P=(v_0,\ldots,v_{m-1})\), put
\[
 A(P)=\frac12\sum_{j<m}\det(v_j,v_{j+1}),\qquad
 e_j=v_{j+1}-v_j,
\]
with cyclic subscripts.  Let rational \(u_j\ge0\) satisfy
\(\lVert e_j\rVert^2\le u_j^2\), and put \(S=\sum_j u_j\).
The previously developed finite theorem is
\begin{equation}\label{eq:plane-finite}
 4\pi_k |A(P)|\le S^2+2^{-k},
 \qquad 0\le\pi-\pi_k<2^{-k},
\end{equation}
where \(\pi_k\) is a specified primitive recursive rational function.
After clearing denominators, this is a universal assertion about
primitive recursive integer calculations.  No simplicity, orientation,
convexity, or nondegeneracy hypothesis is necessary.

Here is the finite reason for the sharp constant.  Let \(N=2^R\ge4\),
\(\eta=c+is\) be the positive half-angle root corresponding to
\(\pi/N\), and \(\zeta=\eta^2\).  These are constructed by the finite
recursion
\[
 c_1=0,\quad s_1=1,\qquad
 c_{r+1}=\sqrt{(1+c_r)/2},\quad
 s_{r+1}=\sqrt{(1-c_r)/2}.
\]
All square roots are positive roots in finite coded quadratic towers.
Their signs and order laws are proved arithmetically in the companion
paper; they are not appeals to an unspecified real-number model.
Write \(x_j+iy_j=\eta^j\).  The upper-half-circle sign table
\[
 y_0=y_N=0,\qquad y_j>0\quad(0<j<N)
\]
has a finite inductive proof: the even powers at the next stage are
the old powers, while each odd power is a positive multiple of the
sum of two adjacent old powers.

For any \(N\)-vertex polygon with complex coordinates \(z_j\), set
\[
 a_h=N^{-1}\sum_{j<N}z_j\zeta^{-jh},\qquad
 E=\sum_{j<N}|z_{j+1}-z_j|^2.
\]
Finite orthogonality gives
\[
 E=4N\sum_{h<N}|a_h|^2 y_h^2,
 \qquad A=N\sum_{h<N}|a_h|^2x_hy_h.
\]
For dyadic \(N\), orthogonality itself is proved by pairing opposite
terms after removing the power of two dividing the frequency.  Thus
neither convergence of a Fourier series nor a geometric-series limit
is involved.  The identity \(cy_h-sx_h=y_{h-1}\) now gives
\begin{equation}\label{eq:plane-dft-positive}
 cE-4sA=4N\sum_{h=1}^{N-1}|a_h|^2y_hy_{h-1}\ge0.
\end{equation}
Conjugating all coordinates reverses the area, and \(c\le1\), so
\(4Ns|A|\le NE\).

To apply this to an arbitrary rational polygon, choose a dyadic
\(N>m\), put \(T=N-m\), and replace edge \(j\) by
\[
 k_j=1+\left\lfloor Tu_j/S\right\rfloor
\]
equal segments; the case \(S=0\) is immediate.  The number of
resulting edges is at most \(N\), so zero-length edges complete the
list to length \(N\).  Bilinearity of the determinant proves that
subdivision preserves area.  Moreover,
\[
 E=\sum_j\frac{\lVert e_j\rVert^2}{k_j}
 \le\sum_j\frac{S u_j}{T}=\frac{S^2}{N-m}.
\]
The inequality includes \(u_j=0\) without dividing by it.
If \(B\ge\max(1,S)\) and
\[
 R\ge k+4,\qquad N\ge m(1+2^kB^2),
\]
the preceding finite identity yields
\[
 4Ns_R|A(P)|\le\frac{N}{N-m}S^2\le S^2+2^{-k}.
\]
The explicit rational lower approximation satisfies
\(\pi_k\le2^{k+4}s_{k+4}\le Ns_R\), proving
\eqref{eq:plane-finite}.  This argument was independently checked
for the present report.

\subsection{Represented boundaries and the equality case}
For closed polygonal maps on a common subdivision, the elementary
determinant identity gives
\begin{equation}\label{eq:plane-area-continuity}
 |A(P)-A(Q)|\le\frac{\|P-Q\|_\infty}{2}
                 \bigl(L(P)+L(Q)\bigr).
\end{equation}
Indeed, if \(d_j=p_j-q_j\), the doubled area difference is
\[
 \sum_j\det(d_j,p_{j+1}-p_j)
 +\sum_j\det(d_{j+1},q_{j+1}-q_j),
\]
and \(|\det(v,w)|\le\|v\|\|w\|\) follows by expanding a
sum-of-squares identity.  Consequently a primitive recursive family
of polygonal approximations, with explicit uniform, area, and length
moduli, satisfies
\[
                 4\pi |A|\le L^2.
\]
For a Jordan-domain presentation, positive enclosed area is supplied
by compatible simple polygonal domain approximations; its name agrees
with the signed-area name after fixing orientation.  This representation
condition is not inferred from bare pointwise injectivity.

The perimeter need not even be a Cauchy real.  With an explicit
variation bound and a uniform-continuity modulus, it can be represented
by the lower cut generated by all inscribed polygon lengths.  A strict
rational lower bound \(q<4\pi|A|\), supplied with a positive rational
margin, determines a sufficiently accurate inscribed polygon by
\eqref{eq:plane-area-continuity}.  Its finite inequality then supplies
a rational lower bound greater than \(q\) for its squared perimeter.
The full margin budget and the conversion of approximating vertices
back to a genuinely inscribed polygon are provided in the companion
paper.  This avoids assuming that the length of every computable
rectifiable curve is computable.

For completeness, the key equality argument is quantitative.  Put
\(q_R=Ns/c\).  The nonzero spectral weights in
\eqref{eq:plane-dft-positive} are
\[
 \mu_h=\frac{4N^2}{c}y_hy_{h-1}\quad(2\le h<N).
\]
They satisfy the exact finite telescope
\[
 \sum_{h=2}^{N-1}\frac1{\mu_h}
 =\frac{c^2}{2N^2s^2}\le\frac18.
\]
To verify it, sum
\(x_{h-1}/y_{h-1}-x_h/y_h=s/(y_hy_{h-1})\)
and use the endpoint values \(c/s\) and \(-c/s\).
Finite weighted Cauchy--Schwarz therefore proves
\begin{equation}\label{eq:plane-finite-stability}
 \max_j|z_j-a_0-a_1\zeta^j|^2
 \le\frac18\bigl(NE-4q_RA\bigr).
\end{equation}
The two unpenalized modes are exactly translation and a positively
oriented regular polygon.

For a closed \(L\)-Lipschitz presentation
\(\gamma:[0,1]\to\mathbb C\), set
\(\Delta=L^2-4\pi A\).  Uniform samples, the area estimate above,
and \(0\le q_R-\pi\le64/N^2\) give
\[
 NE-4q_RA(P_N)\le\Delta+24L^2/N.
\]
The discrete means
\[
 b_N=N^{-1}\sum_j\gamma(j/N),\qquad
 a_N=N^{-1}\sum_j\gamma(j/N)\overline{U(j/N)}
\]
have explicit Cauchy moduli, where \(U\) is the unit-circle
parametrization obtained from the dyadic polygons.  Their limits
\(b,a\) satisfy
\begin{equation}\label{eq:plane-uniform-stability}
 \sup_t|\gamma(t)-b-aU(t)|
 \le\sqrt{\Delta/8+3L^2/N}+8L/N.
\end{equation}
For example, after translating \(\gamma(0)=0\), one has
\(|\gamma|\le L/2\), the functions \(\gamma\) and
\(\gamma\overline U\) have Lipschitz bounds \(L\) and \(5L\),
and dyadic left-endpoint means have error at most half the Lipschitz
bound divided by \(N\).  These facts give precisely the displayed
error budget.  Thus this is a finite quantitative argument, not a
compactness argument about a sequence of nearly circular curves.

If \(\Delta=0\), a dyadic
\[
 N\ge\max(4,48B^2\varepsilon^{-2},16B\varepsilon^{-1}),\quad B\ge L,
\]
makes the error smaller than \(\varepsilon\).  Hence
\(\gamma(t)=b+aU(t)\) in represented equality.  A positive
separation of two boundary points forces \(|a|>0\).
For a general polygonal boundary presentation, apply this argument
to each polygon using rational upper bounds for its edge lengths.
When its deficits tend to zero, the resulting circles converge
without choosing a subsequence: their centers and radii are recovered
from the four coordinate extrema of the limiting boundary.  Thus
equality forces a circle for precisely the effective boundary
representations stated above.  Conversely the dyadic regular polygons
give \(L=2\pi r\), \(A=\pi r^2\) for a circle, so equality holds.

\subsection{Quotients, the two optimization problems, and Dido's problem}
If a positive lower bound for \(L\) is supplied, division is a
primitive recursive operation on the represented names, and
\[
 Q=\frac{4\pi A}{L^2}\le1.
\]
If \(A\) is positively separated from zero, the equivalent ratio
inequality is \(L^2/A\ge4\pi\).  The positivity data avoid an
unbounded search for a nonzero denominator.

A homothety by \(t\ge0\) multiplies the shoelace area by \(t^2\)
and every edge length by \(t\); the same identities hold for the
represented limits.  Therefore a disk attains the maximum area at
fixed positive perimeter and the minimum perimeter at fixed positive
area.  The bounds and their attainment prove these extremal assertions
directly; no separate assertion that an optimizer exists is required.

For Dido's problem, suppose a represented arc of length \(\ell\)
has its endpoints on a line, lies on one side of that line, and,
together with the segment joining its endpoints, presents a domain
of area \(A\).  Reflect the arc across the line.  The resulting
closed boundary has length \(2\ell\) and area \(2A\); the common
straight segment is internal and cancels in the signed-area sum.
Hence
\[
                       A\le\frac{\ell^2}{2\pi}.
\]
Equality implies that the doubled boundary is a circle.  Reflection
invariance puts its center on the given line, and the original arc
is a semicircle.  A semicircle attains equality.  Here the fixed
length is the curved part of the boundary; the endpoints are free
on the line, as in the usual Dido formulation.

\section{Regular polygons: exact quotients and an explicit limit}
\label{sec:regular-new}
For a regular \(n\)-gon, \(n\ge3\), with circumradius \(r>0\), put
\(\alpha=\pi/n\).  Splitting into \(n\) isosceles triangles gives
\[
 L=2nr\sin\alpha,\qquad A=nr^2\sin\alpha\cos\alpha.
\]
Consequently
\begin{equation}\label{eq:regular-quotient}
             Q_n=\frac{\pi}{n\tan(\pi/n)}.
\end{equation}
The cancellation is legitimate because \(0<\alpha<\pi/2\).
For example \(Q_3=\pi/(3\sqrt3)\) and \(Q_4=\pi/4\).

These formulas require only elementary computable trigonometry.  One
may evaluate sine and cosine by their finite Taylor sums on
\([-4,4]\), or use the dyadic unit-circle name from the companion
paper.  For the Taylor route, once the degree exceeds eight the
absolute terms decrease by a factor at most \(1/2\).  The tail after
degree \(M\) is therefore at most \(2^{10-M}\).  Taking degree
\(k+16\), with the argument known within \(2^{-k-16}\), gives a
uniform error below \(2^{-k}\): the derivative sum is bounded by
\(\sum_{j\ge0}4^j/j!<64\).  All assertions are bounds on finite
sums plus a geometric tail estimate.  Addition identities, the
orthogonality integrals used below, and differentiation formulas are
proved by the corresponding finite polynomial identities with these
error bounds; their use does not require an axiom of analytic
completeness.

For \(0<x<\pi/2\),
\[
 \sin x-x\cos x=\int_0^x t\sin t\,dt>0,
\]
so \(Q_n<1\).  The positivity can be given a rational margin:
on \([0,4/3]\), \(\sin t\ge t-t^3/6\ge t/2\), and hence
\(\sin x-x\cos x\ge x^3/6\).
The inequalities \(\sin x\le x\) and
\(\cos x\ge1-x^2/2\) also give
\begin{equation}\label{eq:regular-rate}
             0<1-Q_n\le\frac{\pi^2}{2n^2}\le\frac8{n^2}.
\end{equation}
In particular \(n\ge2^{k+4}\) is an explicit primitive recursive
choice ensuring \(|Q_n-1|\le2^{-k}\).
Moreover \(Q_n\) increases strictly with \(n\):
\[
 \frac{d}{dx}(x\cot x)
   =\frac{\sin x\cos x-x}{\sin^2x}<0,
 \qquad
 x-\sin x\cos x=\int_0^x2\sin^2t\,dt>0.
\]
Each integral inequality here has a precision version proved by
finite Riemann sums and the explicit Taylor bounds.  We use these
identities to prove the elementary statements, not as unexamined
appeals to general real analysis.

\section{The Wigner-caustic improvement as a finite sum of squares}
\label{sec:wigner-new}
The refined planar theorem on the requested page has a particularly
transparent arithmetic form.  We first prove it for a finite
trigonometric support polynomial, then give all approximation data
needed to extend it to effectively presented smooth convex curves.
The coefficient calculation is also the central idea in
Zwierzy\'nski's original paper \cite{zwierzynski}; the treatment below isolates its
finite rational content and its error budget.

\subsection{Support data and the geometric quantities}
Let
\begin{equation}\label{eq:support-polynomial}
 p(\theta)=a_0+\sum_{n=1}^m
                 (a_n\cos n\theta+b_n\sin n\theta),
 \qquad a_n,b_n\in\mathbb Q.
\end{equation}
Write \(u(\theta)=(\cos\theta,\sin\theta)\),
\(v(\theta)=(-\sin\theta,\cos\theta)\), and define
\[
 \gamma(\theta)=p(\theta)u(\theta)+p'(\theta)v(\theta),
 \qquad \rho=p+p''.
\]
For the geometric interpretation assume a supplied positive rational
\(r\) with \(\rho(\theta)\ge r\).  This condition can itself have a
finite certificate.  For example, the sufficient coefficient
condition
\[
 a_0-\sum_{n=2}^m(n^2-1)(|a_n|+|b_n|)\ge r
\]
is rational.  More generally, a rational mesh, interval evaluations,
and a derivative bound for \(\rho\) certify any stated positive
margin.  The theorem does not assert that strict positivity of an
arbitrary represented real function is decidable.

Differentiating the finite expression gives \(\gamma'=\rho v\).
The curve is regular, and it is the positively oriented boundary
of a strictly convex body.  To see this directly, fix \(\theta\).
For \(\theta<\phi\le\theta+\pi\),
\[
 (\gamma(\phi)-\gamma(\theta))\mathbin{\cdot}u(\theta)
   =-\int_\theta^\phi\rho(t)\sin(t-\theta)\,dt
   \le-r\bigl(1-\cos(\phi-\theta)\bigr)<0.
\]
For \(\theta+\pi\le\phi<\theta+2\pi\), integrate backwards
from \(\theta+2\pi\) to obtain the same estimate.  Thus each
tangent line is a supporting line and touches at just its stated
point.  The curve is injective modulo its period, with an explicit
separation bound away from equal parameters.  Inscribed polygons
give the enclosed-area name without requiring an existence theorem
for maximizing convex bodies.

Its length and area are
\begin{equation}\label{eq:support-LA}
 L=\int_0^{2\pi}\rho=\int_0^{2\pi}p,
 \qquad
 A=\frac12\int_0^{2\pi}\det(\gamma,\gamma')
  =\frac12\int_0^{2\pi}(p^2-p'^2).
\end{equation}
For the last equality, expand \(\det(pu+p'v,\rho v)=p\rho\)
and integrate \(pp''\) by parts.  All boundary terms cancel by
periodicity.  These are finite trigonometric calculations; they
also follow by taking polygonal shoelace sums with their explicit
uniform error.  Finite orthogonality now gives
\begin{equation}\label{eq:support-fourier-LA}
 L=2\pi a_0,\qquad
 A=\pi a_0^2-\frac\pi2
       \sum_{n=2}^m(n^2-1)(a_n^2+b_n^2).
\end{equation}
Indeed all mixed-frequency terms integrate to zero,
\(\int\sin^2n\theta=\int\cos^2n\theta=\pi\), and the derivative
multiplies the squared contribution of frequency \(n\) by \(n^2\).

\subsection{The double-cover factor and the exact deficit}
Opposite normal directions give the parallel-tangent pair
\(\gamma(\theta),\gamma(\theta+\pi)\).  Its midpoint is
\[
 w(\theta)=\frac{\gamma(\theta)+\gamma(\theta+\pi)}2
          =f(\theta)u(\theta)+f'(\theta)v(\theta),
 \qquad f(\theta)=\frac{p(\theta)-p(\theta+\pi)}2.
\]
The function \(f\) contains exactly the odd harmonics of \(p\),
and \(w(\theta+\pi)=w(\theta)\).  Thus \([0,2\pi]\) traces
the midpoint curve twice.  The one-pass oriented area appearing in
the Wigner theorem is therefore
\begin{equation}\label{eq:wigner-area}
 \widetilde A_W
  =\frac14\int_0^{2\pi}\det(w,w')
  =\frac14\int_0^{2\pi}(f^2-f'^2)
  =-\frac\pi4
       \sum_{\substack{3\le n\le m\\n\ \text{odd}}}
        (n^2-1)(a_n^2+b_n^2)\le0.
\end{equation}
The factor \(1/4\), rather than \(1/2\), is essential.
Possible cusps or self-intersections cause no ambiguity: this is
signed parametric area, not the unsigned area of the caustic's image.
The original paper likewise counts the double cover in this way.

Subtracting \eqref{eq:wigner-area} from
\eqref{eq:support-fourier-LA} gives the finite identity
\begin{equation}\label{eq:wigner-sos}
 \boxed{L^2-4\pi A-8\pi|\widetilde A_W|
   =2\pi^2
       \sum_{\substack{2\le n\le m\\n\ \text{even}}}
        (n^2-1)(a_n^2+b_n^2)\ge0.}
\end{equation}
In particular the improved inequality is proved, and the ordinary
planar inequality follows by dropping its nonnegative extra term.
The arithmetic heart of \eqref{eq:wigner-sos}, after division by
the known positive constant \(2\pi^2\), is simply nonnegativity
of a finite rational sum of weighted squares.  The displayed
geometric constants are not hiding any nonconstructive step.

\subsection{Equality and a useful quantitative form}
The width in normal direction \(\theta\) is
\[
 W(\theta)=p(\theta)+p(\theta+\pi)
 =2a_0+2\sum_{\substack{2\le n\le m\\n\ \text{even}}}
                   (a_n\cos n\theta+b_n\sin n\theta).
\]
For finite polynomials the following equivalences are immediate and
constructive:
\[
 \begin{split}
 L^2=4\pi A+8\pi|\widetilde A_W|
 &\ \Longleftrightarrow\ a_n=b_n=0\text{ for each even }n\ge2\\
 &\ \Longleftrightarrow\ W(\theta)\text{ is constant}.
 \end{split}
\]
The forward direction uses that a finite sum of nonnegative rational
terms can vanish only if every term vanishes.  The converse direction
can either read the formula or multiply the constant-width identity
by \(\cos n\theta\) and \(\sin n\theta\) and use orthogonality.
The permitted odd harmonics explain why equality in the improved
inequality is much less restrictive than circularity.

There is also a direct estimate that makes the equality argument
survive effective approximation.  Put
\(D_W=L^2-4\pi A-8\pi|\widetilde A_W|\).  Weighted
Cauchy--Schwarz and the finite telescope
\[
 \sum_{j=1}^{M}\frac1{(2j)^2-1}
 =\frac12\left(1-\frac1{2M+1}\right)\le\frac12
\]
give
\begin{equation}\label{eq:width-stability}
       \sup_\theta|W(\theta)-2a_0|^2\le\frac{D_W}{\pi^2}.
\end{equation}
Indeed, for each frequency use
\(|a_n\cos n\theta+b_n\sin n\theta|
 \le\sqrt{a_n^2+b_n^2}\), then apply Cauchy--Schwarz to the
weights \(n^2-1\).  This proves near-constant width from a small
improved deficit with no selection of a minimizing shape.

Similarly the ordinary deficit
\(D=L^2-4\pi A\) satisfies
\[
 \sup_\theta
 |p(\theta)-a_0-a_1\cos\theta-b_1\sin\theta|^2
 \le\frac{3D}{8\pi^2},
\]
because
\(\sum_{n=2}^{M}(n^2-1)^{-1}\le3/4\).
Thus ordinary equality forces
\(p=a_0+a_1\cos\theta+b_1\sin\theta\), and substituting in
\(\gamma=pu+p'v\) gives the circle with center \((a_1,b_1)\)
and radius \(a_0\).  The positive curvature-radius margin implies
\(a_0\ge r>0\).

\subsection{Extension to effective smooth convex support functions}
Here is an explicit representation under which the entire improved
theorem, including its equality characterization, is proved in the
same precision sense as the planar theorem.  The input consists of
a fixed primitive recursive family \(p_j\) of rational
trigonometric polynomials, a bound \(B\ge1\), and a primitive
recursive modulus with
\[
 \|p_j-p\|_\infty,
 \|p'_j-p'\|_\infty,
 \|p''_j-p''\|_\infty\le2^{-j},
\]
after reindexing by that modulus, and with all the displayed
functions bounded in absolute value by \(B\).  In addition supply
a rational \(r>0\) and the inequality \(p+p''\ge r\).
After a prescribed finite initial index, the approximants have
\(p_j+p_j''\ge r/2\), so they represent convex curves as above.
This directly covers ovals with the stated effective normal-angle
support data.  Extracting those data from an arbitrary smooth convex
parametrization, particularly one with vanishing curvature, is an
additional representation problem and is not claimed here.
This is a first-order schema for fixed primitive recursive names;
it is not a universal evaluator of arbitrary primitive recursive
programs.

If two such support data differ in \(C^1\) norm by at most
\(\delta\), with the bound \(B\), direct use of
\(|x^2-y^2|\le2B|x-y|\) gives
\begin{equation}\label{eq:wigner-error-budget}
 |L_p-L_q|\le8\delta,\qquad
 |A_p-A_q|\le16B\delta,\qquad
 |\widetilde A_{W,p}-\widetilde A_{W,q}|\le8B\delta.
\end{equation}
Here \(\pi\le4\); the factor for the last term is half the
area factor because the midpoint curve has the one-pass convention
of \eqref{eq:wigner-area}.  Since \(|L_p|,|L_q|\le8B\),
\[
                         |D_{W,p}-D_{W,q}|\le640B\delta.
\]
Choosing \(\delta\le2^{-k}/(1024B)\), and increasing the rational
precision for the finite coefficient calculations, gives the
requested precision inequality.  Thus \eqref{eq:wigner-sos}
passes to the represented limit with an explicit primitive recursive
error budget.  So does \eqref{eq:width-stability}, since the widths
and their means have the same supplied convergence modulus.
If \(D_W=0\), it follows directly that \(W\) is constant.

Conversely, put \(q=(p(\theta)+p(\theta+\pi))/2\).  The
finite-polynomial formula for the improved deficit can also be written
\[
 D_W=2\pi\int_0^{2\pi}
                 \bigl(q'^2-(q-a_0)^2\bigr)\,d\theta,
 \qquad a_0=(2\pi)^{-1}\int_0^{2\pi}p.
\]
It remains valid by \eqref{eq:wigner-error-budget}.  If the width
is constant, \(q=a_0\) and \(q'=0\), so \(D_W=0\).
Derivative uniqueness here is effective: for a function whose
derivative has modulus \(\omega\), its difference quotient over
an explicitly small positive \(h\) differs from its derivative
by the corresponding modulus error.  For a constant function the
difference quotient is exactly zero.  This argument uses no
principle that a real number is either zero or positively separated
from zero.

The representation is not confined to preassigned Fourier series.
For a primitive-recursively evaluable periodic \(C^2\) function
with named derivatives, a common uniform-continuity modulus and
a bound \(B\) for \(p,p',p''\), a suitable polynomial family
can be constructed by the following finite Fej\'er procedure.
Define
\[
 F_N(t)=N^{-1}\left|\sum_{h=0}^{N-1}e^{iht}\right|^2.
\]
This is a nonnegative trigonometric polynomial and has integral
\(2\pi\), by finite orthogonality.  For circular distance
\(|t|\ge\delta\), \(0<\delta\le1\), the finite geometric-sum
identity and \(\sin(|t|/2)\ge|t|/\pi\ge\delta/4\) give
\(F_N(t)\le16/(N\delta^2)\).
If the common modulus gives oscillation at most \(2^{-k-1}\)
within \(\delta=2^{-\omega(k+1)}\), then convolution satisfies
\[
 \left\|\frac1{2\pi}\int_0^{2\pi}F_N(t)p(\,\cdot-t)\,dt-p
          \right\|_{C^2}
 \le2^{-k-1}+\frac{32B}{N\delta^2}.
\]
Take \(N\ge64B\,2^k\delta^{-2}\).
Differentiation commutes with this finite-kernel convolution by
periodic integration by parts; the identity is also obtained by
finite summation by parts with a modulus error tending to zero.
Approximate its finitely many Fourier coefficients rationally.
An error at most \(\tau\) in each coefficient changes the
\(C^2\) norm by at most \(6N^3\tau\).  Taking the Fej\'er
target one bit more accurate and
\(\tau\le2^{-k-1}/(6N^3)\) yields the required rational
polynomial at accuracy \(2^{-k}\).
The Fourier integrals are computed by uniform Riemann sums using
the given modulus, the frequency bound, and the elementary
trigonometric error bound.  Every loop count is specified before
the loop starts.  Thus smoothness with effective moduli supplies
the representation; bare classical smoothness without those data
is not being claimed to do so.

\subsection{Arithmetic and formalization boundary}
For finite rational coefficients introduce the rational sums
\[
 E=\sum_{n=2}^m(n^2-1)(a_n^2+b_n^2),\quad
 O=\sum_{\substack{3\le n\le m\\n\ \text{odd}}}
          (n^2-1)(a_n^2+b_n^2),\quad H=E-O.
\]
The internal finite theorem is \(H\ge0\), together with the
explicit identities \(L^2-4\pi A=2\pi^2E\) and
\(-\widetilde A_W=\pi O/4\), understood through rational
names for \(\pi\).  The partition \(E=O+H\), nonnegativity,
and the zero-sum characterization are proved by induction on the
finite coefficient list.  After coding rational arithmetic, the
invariants are quantifier-free equalities or inequalities between
primitive recursive folds.  The infinite-looking statements above
mean that a requested precision selects the finite index in
\eqref{eq:wigner-error-budget}; no induction over an analytic
predicate is used.

The Lean development should first formalize these finite rational
folds, then the finite trigonometric identities and explicit Riemann
sum estimates, and finally the represented support-function schema.
The finite sum-of-squares certificate and the complete modulus budget
are essential: a theorem about Mathlib's real numbers alone would
not establish derivability in the repository's embedded
\(\mathsf{iPRA}\).
